\printedpage{104}
% head: A. Grothendieck : Dérivateurs, Ch. VII, Édition des Universités de Montpellier II et Paris VII
\msp{page 166}

Thus, the functor $\psi\varphi$ is indeed homotopic to the identity functor. We have won!

Let us pass to $1'$, $2'$) for $\CN$. The assertion $1'$) is proved in exactly the same way as 1) for \underline{Ch}$(X)$, by regarding $\CN$ as fibered over $\mathbf{N}^o$ (which has a final object 0), and using the diagram
\[
\begin{tikzcd}
X = (X_{\mathbf{N}^o})_e \arrow[rr, "i'"] & & X_{\mathbf{N}^o} = X \times \mathbf{N}^o \arrow[rr, "j'"] \arrow[dr] & & \mathcal{C}_{\mathbf{N}^o} \arrow[dl] \\
 & & & \mathbf{N}^o \ , &
\end{tikzcd}
\]
where $j'$ is (cartesian and) a Hot-equivalence on each fiber by the lemma, hence a Hot-equivalence globally, and $i'$ is a Hot-equivalence since $\mathbf{N}^o$ is aspheric (having a final object). Thus $j'i'$ is aspheric, q.e.d.

It remains to prove that
\[ \CN \longrightarrow X \times X \]
is proper and smooth. Here again,
\msp{page 167}
it is exactly the same argument as before, I presume. There is, however, a technical complication. We had taken, to define
\[ \underline{t} : \mathbf{N}^o \longrightarrow \underline{\mathrm{Ty}}, \qquad \underline{t} = (t_n)_n, \]
the choice
\[ t_n = (\Delta_+\Delta_-)^{2n}, \]
or if one prefers
\[ t_n = (\Delta_+\Delta_-)^n, \]
in any case the $t_n$ have initial direction $+$, terminal $-$, and they are reduced types. If one takes $\Delta_+t_n\Delta_-$, one finds a nonreduced type, which is not among the $t_n$. It would be of no use to change the definition of $t_n$ to $t_n = (\Delta_-\Delta_+)^{\text{something}}$, for this time $\Delta_+t_n\Delta_-$ will indeed be reduced, but will not have the right initial and terminal signs. Thus to define $\tilde{c}$ in terms of $c$, as an object of $\CN$, it will be necessary, if $c$ is of type $\tau = t_n$, to take $\tilde{c}$ of type $\Delta_+\Delta_-\tau\Delta_+\Delta_- = t_{n+1}$ (taking $t_{n+1} = (\Delta_+\Delta_-)^{2n}$), and likewise to define $\tilde{c}'$ in terms of $c'$. The relevant definitions of the passage from $(c',w')$ to $(\tilde{c}',\tilde{w}')$ are described in the figure: % [sic] t_{n+1} = (\Delta_+\Delta_-)^{2n} -> t_{n+1} = (\Delta_+\Delta_-)^{2n+2}, exponent inconsistent with t_n = (\Delta_+\Delta_-)^{2n}

\msp{page 168}
% foot: Version 22 décembre 2001

\printedpage{105}
% head: A. Grothendieck : Dérivateurs, Ch. VII, Édition des Universités de Montpellier II et Paris VII
\[
\tilde{w}'\Big\downarrow
\left\{
\overbrace{\underbrace{
\begin{tikzcd}[ampersand replacement=\&]
a' \arrow[r, "\mathrm{id}"] \arrow[d, "\mathrm{id}"'] \& a' \arrow[d, "u"'] \& a' \arrow[l, "\mathrm{id}"'] \arrow[d, "u"'] \arrow[rrrr, dotted, no head, "c'", "\longrightarrow"'] \& {} \arrow[d, "w'"] \& {} \arrow[d, "w'"] \& {} \arrow[d, "w'"] \& b' \arrow[r, "\mathrm{id}"] \arrow[d, "v"'] \& b' \arrow[d, "v"] \& b' \arrow[l, "\mathrm{id}"'] \arrow[d, "\mathrm{id}"] \\
a' \arrow[r, "u"] \& a \& a \arrow[l, "\mathrm{id}"'] \arrow[rrrr, no head, "c", "\longrightarrow"'] \& {} \& {} \& {} \& b \arrow[r, "\mathrm{id}"] \& b \& b'. \arrow[l, "v"']
\end{tikzcd}
}_{\tilde{c}}}^{\tilde{c}'}
\right\}
\tilde{w}'
\]

Thus
\[
\left\{
\begin{array}{lcl ll}
\tilde{c} & = & \xi_u \circ c \circ \eta_v & \text{where} &
\left\{
\begin{array}{lccccc}
\xi_u : & a' & \xrightarrow{u} & a & \xleftarrow{\mathrm{id}} & a \\[1ex]
\eta_v : & b & \xrightarrow{\mathrm{id}} & b & \xleftarrow{v} & b'
\end{array}
\right\}
\text{ type } \Delta_+\Delta_- \\[4ex]
\tilde{c}' & = & a'_{\Delta_+\Delta_-} \circ c' \circ b'_{\Delta_+\Delta_-} & & \\[4ex]
\tilde{w}' & = & \lambda(u) * w * \mu(v) & \text{with} & % [sic] w -> w', the homomorphism of the pair (c',w')
\left\{
\begin{array}{lccc}
\lambda(u) : & a'_{\Delta_+\Delta_-} & \longrightarrow & \xi_u \\[1ex]
\lambda(v) : & b'_{\Delta_+\Delta_-} & \longrightarrow & \eta_v % [sic] \lambda(v) -> \mu(v), the formula above has \mu(v)
\end{array}
\right\}
\text{ as in \texttt{[the]} figure.}
\end{array}
\right.
\]
One thus finds, mutatis mutandis,
\[ \CN(c;u,v) \xrightarrow{\varphi} \CN(a',b')/\tilde{c}, \qquad \varphi(c',w') = (\tilde{c}',\tilde{w}'), \]
and one defines the functor
\[ \CN(c;u,v) \xleftarrow{\psi} \CN(a',b')/\tilde{c} \]
as follows: to $(c'',w'')$ in $\CN(a',b')/\tilde{c}$, thus a path $c''$ of type $t_n = (\Delta_+\Delta_-)^{2n}$ from $a'$ to $b'$ and a homomorphism $w'' : c'' \longrightarrow \tilde{c}$ with fixed endpoints, one associates $(c',w')$ in $\CN(c;u,v)$, with
\[ c' = c'', \qquad w' = \rho_c w'', \]
where
\[ \rho_c : \tilde{c} \longrightarrow c \]
is defined in a similar way as before, by
\msp{page 169}
\[
\begin{tikzcd}
a' \arrow[r, "u"] \arrow[d, "u"'] & a \arrow[d, "\mathrm{id}"] & a \arrow[l, "\mathrm{id}"'] \arrow[d, "\mathrm{id}"] \arrow[rr, zigzag] & {} \arrow[d, "\mathrm{id}"'] & b \arrow[r, "\mathrm{id}"] \arrow[d, "\mathrm{id}"] & b \arrow[d, "\mathrm{id}"] & b' \arrow[l, "v"'] \arrow[d, "v"] & \tilde{c},\ \text{ type } t_{n+1}. \\
a \arrow[r, "\mathrm{id}"] & a & a \arrow[l, "\mathrm{id}"'] \arrow[rr, zigzag, "c"] & {} & b \arrow[r, "\mathrm{id}"] & b & b' \arrow[l, "\mathrm{id}"'] & \varphi^*_{n,n+1}(c).
\end{tikzcd}
\]

One need only check that the endofunctor $\psi\varphi$ of $\CN(c;u,v)$ is homotopic to the identity. But this functor transforms $(c',w')$ into $(\tilde{c}' = a'_{\Delta_+\Delta_-}c'b'_{\Delta_+\Delta_-}, \overline{w'})$, where $\overline{w'}$ is again the composite
% foot: Version 22 décembre 2001
\printedpage{106}
% head: A. Grothendieck : Dérivateurs, Ch. VII, Édition des Universités de Montpellier II et Paris VII
\[
\begin{tikzcd}
\tilde{c}' \arrow[r, "\text{can.}"] \arrow[rr, bend right=20, "\overline{w'}"'] & c' \arrow[r, "w'"] & c.
\end{tikzcd}
\]

We therefore have
\[ \psi\varphi \longrightarrow \mathrm{id}, \]
$\psi\varphi$ is homotopic to the identity functor.

This proves the smoothness of $\CN \longrightarrow X \times X$. For properness, one must prove the smoothness of the case of $X^o$, and of the inverse system $\underline{t}^{\mathrm{opp}} \simeq (t_n^{\mathrm{opp}})_n$. But we have $\underline{t}^{\mathrm{opp}} = \underline{t}$, and it suffices to apply what we have already proved to the case of $X^o$, to see that $\CN \longrightarrow X \times X$ is proper. % [sic] de le cas -> dans le cas, grammatical

Only the lemma remains to be proved. It must be essentially
\msp{page 170}
trivial, and I give up writing out a proof.

I am now going to show ``formally'' that the ``canonical'' factorization of an arrow in Cat $f : X' \longrightarrow X$ as
\[ X' \xrightarrow{i_f} \overline{X'} \xrightarrow{\bar{f}} X, \qquad i_f \text{ Hot-equivalence,} \]
with (we hoped) $\bar{f}$ proper and smooth, starting from a similar factorization for the diagonal arrow $X \xrightarrow{\ \ \mathrm{diag}_X\ \ } X \times X$ as
\[
\begin{tikzcd}
X \arrow[d, "i_X\ \ \text{\footnotesize Hot-equivalence}"] \\
\mathcal{C}_X\ \ \text{``paths''} \arrow[d, "p_X"] \\
X \times X
\end{tikzcd}
\]
is a formal matter. (Cf. page 139 for the notations and \texttt{[the]} problematics.) We look at the diagram
\[
\begin{tikzcd}
 & \mathcal{C}_X \arrow[d, "p_X"] \arrow[dl, "\beta_X"'] \arrow[dr, phantom, "\text{cart.}"] & C(f) \arrow[l, "\mathrm{pr}'_2"'] \arrow[d, "p_{X,f}"] \arrow[dll, bend right=50, "\bar{f}"'] & [-1.5em] \simeq X' \times_X (C_X, \sigma = \mathrm{pr}_1 \circ p_X) \\ % Reading resolved against the supplied PDF; see corrections.csv.
X & X \times X \arrow[l, "\mathrm{pr}_2^X"] \arrow[d, "\mathrm{pr}_1^X"'] \arrow[dr, phantom, "\text{cart.}"] & X' \times X \arrow[l, "f \times \mathrm{id}_X = f'"] \arrow[d, "\mathrm{pr}_1"] & \\
 & X & X' \arrow[l, "f"'] \arrow[uu, bend right=60, dotted, "i_f"'] &
\end{tikzcd}
\]
% foot: Version 22 décembre 2001
\printedpage{107}
% head: A. Grothendieck : Dérivateurs, Ch. VII, Édition des Universités de Montpellier II et Paris VII
(\footnote[40]{NB $C(f)$ is defined as the inverse image over $X'$ of the bundle $C_X$ over $X$ (by $\sigma_X = \mathrm{pr}_1^X \circ p_X$) by $f : X' \longrightarrow X$.}). Since $p_X$ is proper and smooth, i.e. ``perfect'', its inverse image over $X' \times X$, namely $p_{X,f}$, is so too; on the other hand $\mathrm{pr}_2^X \circ f'$ is the
\msp{page 171}
projection $\mathrm{pr}_2$ of $X' \times X$ onto $X$, hence it is also a perfect morphism (since $X' \longrightarrow e$ is perfect, hence the same holds for the composite
\[ (\mathrm{pr}_2^X \underbrace{f') \circ p_{X,f}} = \underbrace{\mathrm{pr}_2^X(p_X}\, \mathrm{pr}'_2) = \beta_X \mathrm{pr}'_2 \overset{\text{def}}{=} \bar{f} \ \ ). \]

I recall that the section $i_f$ of $C(f)$ over $X'$ is defined from the section $i_X$ of $(\mathcal{C}_X,\sigma_X)$ over $X$. It is a Hot-equivalence, for $\sigma_X = \mathrm{pr}_1^X p_X$ is perfect with aspheric fibers, hence universally a Hot-equivalence, hence $\mathrm{pr}_1 p_{X,f} : C(f) \longrightarrow X'$ is also a Hot-equivalence, and any section is therefore one too.
% foot: Version 22 décembre 2001